\section*{Phụ lục: Khai báo sử dụng Generative AI}
\addcontentsline{toc}{section}{Phụ lục: Khai báo sử dụng Generative AI}

Theo yêu cầu minh bạch của môn học, bảng dưới đây tổng hợp việc sử dụng công cụ Generative AI của các thành viên trong Nhóm 09 (cập nhật từ tệp \texttt{docs/AI\_USAGE.md}).

\begin{longtable}{|p{2.3cm}|p{2.1cm}|p{3.0cm}|p{2.5cm}|p{3.5cm}|}
\caption{Bảng khai báo sử dụng Generative AI của Nhóm 09} \\
\hline
\textbf{Thành viên} & \textbf{Công cụ} & \textbf{Dùng cho phần nào} & \textbf{Mức độ đóng góp} & \textbf{Nhóm đã tự chỉnh sửa, kiểm tra gì} \\
\hline
\endfirsthead
\hline
\textbf{Thành viên} & \textbf{Công cụ} & \textbf{Dùng cho phần nào} & \textbf{Mức độ đóng góp} & \textbf{Nhóm đã tự chỉnh sửa, kiểm tra gì} \\
\hline
\endhead
TV1 (Nguyễn Khang Kỳ) & \textit{\_TODO\_} & \textit{\_TODO\_} & \textit{\_TODO\_} & \textit{\_TODO\_} \\ \hline
TV2 (Châu Mạnh Phát) & \textit{\_TODO\_} & \textit{\_TODO\_} & \textit{\_TODO\_} & \textit{\_TODO\_} \\ \hline
TV3 (Hoàng Phạm Tiến Thành) & Claude \& agent lập trình & Gợi ý phân công nhóm, dựng khung Flask MVP, viết prompt dựng giao diện Dashboard từ mockup Figma & Gợi ý ý tưởng, viết nháp, sinh code màn hình Dashboard & Tự chạy thử test/server xác nhận khung code chạy đúng, tự chốt nhãn hiển thị/nút bấm, đối chiếu với mockup gốc \\ \hline
TV4 (Võ Tấn Phát) & ChatGPT, Claude & Gợi ý test case cho Submission 3 và kiểm tra lỗi code & Gợi ý ý tưởng, kiểm tra lỗi & Kiểm tra lại tính phù hợp của test case, tự sửa lỗi code bằng tay \\ \hline
TV5 (Đào Huy Phong) & \textit{\_TODO\_} & \textit{\_TODO\_} & \textit{\_TODO\_} & \textit{\_TODO\_} \\ \hline
\end{longtable}
